\subsubsection*{准静态为什么还不够}

准静态只表示变化足够缓慢，使系统能用一串近似平衡态描述；可逆还要求系统和环境都能在不留下
其他变化的条件下复原。例如，缓慢拖动有摩擦的活塞可以近似准静态，但摩擦仍把机械能转为内能，
单纯反向拖动不能把这些内能全部恢复为原来的机械能。有限温差传热也有同样的问题。
因此，可逆传热须令两侧温差趋于零，可逆体积功须令内外压强差趋于零，并排除摩擦等耗散。
这是一种极限模型，不是要求实际装置在有限时间内严格达到它。

\subsubsection*{熵流与熵产生}

对不交换物质的系统，热量可以从不同温度的环境输入。设第 $j$ 个恒温热库温度为 $T_j$，
在全过程向系统传入的带符号热量为 $Q_j$。热库自身的熵变为 $-Q_j/T_j$。
把系统和这些热库一起考察，第二定律要求总熵变非负，因而
\begin{equation}
 \Delta S_{\rm sys}-\sum_j\frac{Q_j}{T_j}=S_{\rm gen}\ge0.
\end{equation}
这里 $S_{\rm gen}$ 是系统连同有关传热边界的总熵产生；它不是新的能量。
等号对应整个过程可逆。如果环境温度变化，应把 $Q_j/T_j$ 换成相应的积分，不能随意选一个平均温度。
对循环，$\Delta S_{\rm sys}=0$，于是
\begin{equation}
 \sum_j\frac{Q_j}{T_j}\le0.
\end{equation}
这就是以恒温热库分段表示的克劳修斯不等式。循环可以把吸热和放热的熵交换相互抵消，
却不能用负的熵产生去抵消耗散。

例如，两热库间直接传递热量 $Q>0$，高温为 $T_h$、低温为 $T_c<T_h$。
高温热库熵变 $-Q/T_h$，低温热库熵变 $Q/T_c$，故
\[
 S_{\rm gen}=Q\left(\frac1{T_c}-\frac1{T_h}\right)>0.
\]
总能量没有损失，但可以利用的温差减小了。若 $T_h\to T_c$，单位热量引起的熵产生趋于零，
这正是可逆传热必须取无穷小温差的原因。

\subsubsection*{由熵不等式判断热机与制冷机}

在 $T_h>T_c>0$ 的两个热库之间，热机每循环从高温热库取热 $Q_h>0$，向低温热库排热 $Q_c>0$。
工质循环复原，环境总熵变为
\[
 -\frac{Q_h}{T_h}+\frac{Q_c}{T_c}\ge0
 \quad\Longrightarrow\quad
 \frac{Q_c}{Q_h}\ge\frac{T_c}{T_h}.
\]
结合 $W=Q_h-Q_c$，得到
\begin{equation}
 \eta=1-\frac{Q_c}{Q_h}\le1-\frac{T_c}{T_h}.
\end{equation}
这次推导没有使用气体物态方程，所以卡诺效率是任意工质的上限，而不只是理想气体计算的巧合。
若装置的某一部分不可逆，通常有严格小于号。

制冷机的有用效果是从低温热库取出 $Q_c$，外界输入功 $W>0$，向高温热库排出
$Q_h=Q_c+W$。此时热库熵变满足
\[
 \frac{Q_h}{T_h}-\frac{Q_c}{T_c}\ge0
 \quad\Longrightarrow\quad
 W\ge Q_c\frac{T_h-T_c}{T_c}.
\]
制冷系数和供热系数分别为
\begin{equation}
 \varepsilon_{\rm ref}=\frac{Q_c}{W}\le\frac{T_c}{T_h-T_c},\qquad
 \varepsilon_{\rm hp}=\frac{Q_h}{W}=1+\varepsilon_{\rm ref}
 \le\frac{T_h}{T_h-T_c}.
\end{equation}
二者可以大于1，因为它们比较的是搬运热量与输入功，而不是把输入热量转为功的比例。
例如 $T_h=300\,\mathrm K,T_c=270\,\mathrm K$ 时，可逆制冷系数为9；搬走
$900\,\mathrm J$ 的热量至少需要 $100\,\mathrm J$ 的功，向高温环境排热为
$1000\,\mathrm J$。由 $1000/300=900/270$ 可同时核验熵平衡。
